\section{A determinant obstruction for bounded inducing images}
\label{sec:bounded-inducing-determinants}

\paragraph{Status, provenance and scope.}
These are HD-PD-01--02, with question row Q2cy reserved by root.
The simultaneous-prime obstruction and its integral strengthening
originated with root.  The polarized descent, regular-trace
matrix-order proof and homothetic-lattice proof were developed
jointly with /root/imprimitive.  /root/g34\_eightfolds has
independently read the complete message derivations and then the whole
saved proof. Author, root and separate peer frozen-file readings all pass;
the shared derivation is not represented as independent authorship.
HD-PD-02 is unconditional for the stated bounded inducing-image
class.  It does not assert that a previously enumerated structural
class exhausts all finite-group sources, and does not obstruct
arbitrary algebraic cycles.

\subsection*{HD-PD-01. Polarized Frobenius reciprocity}

Fix an imaginary quadratic field $K$ and an element
$\eta\in K^*$ with $\bar\eta=-\eta$.
Use the same $\eta$ in all Hermitian normalizations.
Suppose a smooth complex curve $C$ has a finite $G$-action.
Let $\rho_K$ be an absolutely irreducible $K$-model with exact
character field $K$, and let $B$ be the six-dimensional primitive
factor selected by a rational Rosati-fixed rank-one projector in
its $\mathbf Q[G]$ component.  Its polarization is the restriction
of the canonical Jacobian polarization, and its $K$-signature
is $(3,3)$.  Write its signed determinant as
\[
 \delta_B=(-1)^3\det(h_B)
    \quad\hbox{in}\quad
 \mathbf Q^*/\operatorname{Nm}_{K/\mathbf Q}K^*.
 \tag{PD.1}
\]
Assume
$\rho_{\mathbf C}=\operatorname{Ind}_H^G\tau$ with $\tau$
irreducible, put $N=\ker\tau$, $J=H/N$, and $Y=C/N$.
The quotient map is $\pi:C\to Y$.

\begin{proposition}[HD-PD-01]
Let $E$ be a CM cyclotomic field containing $K$ and splitting
$J$, and put $F=E^+$.  The rank-six $\tau$-multiplicity Hermitian
form in $H_1(Y,E)$ has signed determinant equal to the image of
$\delta_B$ in $F^*/\operatorname{Nm}_{E/F}E^*$.
No generic endomorphism-algebra or period-rank assumption is
required.  The field $E$ need not split the original group $G$.
\end{proposition}

\begin{proof}
Choose a model $\tau_E$.  The two representations
$\operatorname{Ind}_H^G\tau_E$ and $\rho_K\otimes_KE$ have
isomorphic complexifications, which are absolutely irreducible.
They are therefore isomorphic over $E$.
We use homology and left deck action consistently.  The extension
of the alternating intersection pairing $I_C$ gives the Hermitian
pairing
\[
                  H_C(u,v)=\eta^{-1}I_C(u,\bar v).
 \tag{PD.2}
\]
The bar on vectors conjugates the coefficient field; it is not
a new action on the underlying integral homology.
On the $\rho$-isotypic summand
$V_\rho\otimes_E M_\rho$, invariance and Schur's lemma give
\[
                         H_C=h_\rho\otimes b_\rho,
 \tag{PD.3}
\]
where $h_\rho$ is a fixed positive invariant $K$-Hermitian
coefficient form.  The multiplicity space has dimension six.
The rational primitive projector selecting $B$ selects a
$K$-line of coefficient norm $a\in\mathbf Q_{>0}$.
Thus the natural Hermitian form of $B$, under the corresponding
identification, is $a b_\rho$.  Its determinant differs by
$a^6=\operatorname{Nm}_{K/\mathbf Q}(a^3)$.

Frobenius reciprocity gives
\[
 \dim_E\operatorname{Hom}_H
          (\tau_E,\operatorname{Res}\rho_E)=1.
 \tag{PD.4}
\]
It also excludes $\tau_E$ from every other $G$-constituent:
any such constituent would receive a nonzero map from
$\operatorname{Ind}\tau_E=\rho_E$.
For a nonzero intertwiner $i:V_\tau\to V_\rho$ and a fixed
positive invariant coefficient form $h_\tau$, Schur's lemma gives
\[
                       i^*h_\rho=c h_\tau,\qquad c\in F_{>0}.
 \tag{PD.5}
\]
Thus the entire $\tau$-multiplicity form inside $H_1(C,E)$ is
$c b_\rho$.  Jacobian pullback $\pi^*:J(Y)\to J(C)$ induces
the homological transfer to $N$-invariants.  It multiplies the
intersection pairing by $|N|$.  Hence, after the actual
multiplicity-space identification,
\[
                        b_{\tau,Y}=\frac{c}{|N|}\,b_\rho.
 \tag{PD.6}
\]
Changing a coefficient form, choosing another self-adjoint
primitive coefficient line, or changing the anti-invariant
normalization multiplies this form by an element of $F^*$.
Every such rank-six determinant factor is a norm:
\[
                 s^6=\operatorname{Nm}_{E/F}(s^3)
                 \qquad(s\in F^*).
 \tag{PD.7}
\]
This proves the assertion.  In particular, no assumption that
$p$ is prime to the original group order, and no generic
endomorphism calculation, is hidden in this comparison.
\end{proof}

\paragraph{The weaker finite-catalogue form.}
Suppose a finite set of inducing topological entries exhausts a
chosen source class.  Their multiplicity pairings give finitely
many fixed determinant classes.  Pass to a common cyclotomic
field containing their coefficient fields and $K$, and choose
representatives.  A prime congruent to $-1$ modulo its conductor,
outside the finite supports of those representatives, has
Frobenius equal to complex conjugation.  It splits completely
in the real fixed field and is inert in the CM quadratic
extension at every place above it.  Every representative has
valuation zero, whereas the target class $p$ has valuation one;
every local norm has even valuation.  HD-PD-01 excludes this
target class for all entries at once.  Dirichlet's theorem gives
infinitely many such primes \cite[Section 1, p.~63]{PD:Harrison2009}.
This is conditional on the stated exhaustiveness hypothesis.
The next theorem removes the representative list and support
calculation from the bounded-order case.

\subsection*{HD-PD-02. The integral obstruction depends only on the inducing order}

\begin{theorem}[HD-PD-02]
For any integer $Q\ge1$, put
\[
            M_Q=\operatorname{lcm}(8,|\operatorname{disc}K|,
                                      1,2,\ldots,Q).
 \tag{PD.8}
\]
For every prime $p\equiv-1\pmod {M_Q}$ and every naturally
polarized primitive sixfold in HD-PD-01 with $|J|\le Q$,
\[
                             v_p(\delta_B)\equiv0\pmod2.
 \tag{PD.9}
\]
Consequently the signed target class $p$ is impossible.
The original group, induction index, curve genus, quotient
genus and branch ramification are unrestricted.

There are infinitely many such pairwise distinct target
determinant classes.  A terminating prime search constructs
them from $K,Q$ without GRH or enumeration of branch tuples.
The exclusion also applies under $K$-isogeny to a target with
$\operatorname{End}^0=K$, even when its specified polarization
was not initially required to match the natural source polarization.
\end{theorem}

\begin{proof}
Set $M=M_Q$ and $E=\mathbf Q(\zeta_M)$.
The field contains $K$ and splits every $J$ of order at most $Q$,
since it contains all roots of unity of exponent dividing $|J|$.
This is the usual cyclotomic splitting-field consequence of
Brauer induction; see
\cite[Corollary 4.6.12 and its proof]{OG:Zheng2025}
and the previously checked Brauer input.

For $p\equiv-1\pmod M$, cyclotomic Frobenius is complex
conjugation.  Thus every completion $E_w/\mathbf Q_p$ is
unramified quadratic, its fixed field is $\mathbf Q_p$, and
$K_p=E_w$.  These Frobenius and tower assertions are
\cite[8.15--8.18, pp.~142--143]{PD:MilneANT}.
In particular $p$ is odd, $p\nmid |J|$, and
$\eta=\sqrt{\operatorname{disc}K}$ is an anti-invariant $p$-unit.
Let $\mathcal O$ be the integers of $E_w$.
The local field $E_w$ contains all $\exp(J)$-th roots of unity
and splits $J$.  No assertion that it splits $G$ is needed.

\paragraph{The integral group algebra is a product of matrix orders.}
For every irreducible $E_w[J]$-module $V_\alpha$, choose a
$J$-stable $\mathcal O$-lattice $\Lambda_\alpha$.  For instance,
take the $\mathcal O$-span of the finite $J$-orbit of a basis.
Under the split group-algebra isomorphism set
\[
 \mathcal A'=\prod_\alpha
       \operatorname{End}_{\mathcal O}(\Lambda_\alpha)
       \ \supseteq\ \mathcal O[J].
 \tag{PD.10}
\]
Consider the regular trace pairing
$T(x,y)=\operatorname{Tr}(L_{xy}|E_w[J])$.
On the group basis,
\[
             T(g,h)=|J|\,\mathbf 1_{\{gh=1\}}.
 \tag{PD.11}
\]
Since $|J|$ is a unit, $\mathcal O[J]$ is its own trace-dual.
For $x\in\mathcal A'$ and $y\in\mathcal O[J]$,
\[
 T(x,y)=\sum_\alpha \dim(V_\alpha)
                       \operatorname{tr}(x_\alpha y_\alpha)
                       \ \in\ \mathcal O.
 \tag{PD.12}
\]
Hence $x$ belongs to the trace-dual of $\mathcal O[J]$.
This proves
\[
                  \mathcal O[J]=\mathcal A'.
 \tag{PD.13}
\]
In particular the integral matrix-order assertion has been
proved directly, without assuming a maximal-order classification.

\paragraph{A unimodular coefficient form.}
The involution
$*:\sum a_g g\mapsto\sum\bar a_g g^{-1}$
fixes every primitive character central idempotent, because
$\overline{\chi_\alpha(g)}=\chi_\alpha(g^{-1})$.
It preserves the matrix orders in (PD.13).
Choose an invariant nondegenerate Hermitian coefficient form
$h_\alpha$; the positive averaged form of a global cyclotomic
model supplies one.  Its adjoint involution is $*$.
Consequently
\[
 \operatorname{End}_{\mathcal O}(\Lambda_\alpha^\#)
       =\operatorname{End}_{\mathcal O}(\Lambda_\alpha).
 \tag{PD.14}
\]
Here $\#$ denotes the dual lattice for $h_\alpha$.
Two full lattices with the same full matrix endomorphism
order are homothetic: in a basis of the first, the diagonal
matrix units split the second into coordinate fractional
ideals, and the off-diagonal units force all ideals equal.
Since $\mathcal O$ is a discrete valuation ring,
$\Lambda_\alpha^\#=p^k\Lambda_\alpha$ for an integer $k$.
Rescaling $h_\alpha$ by $p^k\in\mathbf Q_p$ makes
$\Lambda_\alpha$ self-dual.

\paragraph{A unimodular multiplicity form.}
The module
$U=H_1(Y,\mathbf Z_p)\otimes_{\mathbf Z_p}\mathcal O$
has the perfect Hermitian pairing
\[
                    H(u,v)=\eta^{-1}I_Y(u,\bar v).
 \tag{PD.15}
\]
The integral intersection form $I_Y$ is unimodular and
$\eta$ is a unit.  Deck transformations act unitarily.
Every self-adjoint central summand $U_\alpha=e_\alpha U$
is therefore perfect.  Morita decomposition for the full
matrix order gives
\[
 U_\alpha=\Lambda_\alpha\otimes_{\mathcal O}W_\alpha,
          \qquad H_\alpha=h_\alpha\otimes b_\alpha.
 \tag{PD.16}
\]
The second equality follows from the adjoint compatibility
with the full matrix algebra.  The modules $W_\alpha$ are
free over $\mathcal O$.

For the relevant $\tau$ constituent, Frobenius reciprocity
and the dimension of $B$ give $\operatorname{rank}W_\tau=6$.
The form $b_\tau$ is integral: a unimodular matrix of
$h_\tau$ has some unit entry, and each coefficient of $b_\tau$
is obtained from the corresponding integral coefficient of
$H_\tau$ by division by that entry.  Moreover
\[
        \det H_\tau=(\det h_\tau)^6
                            (\det b_\tau)^{\dim\tau}.
 \tag{PD.17}
\]
Both $H_\tau$ and $h_\tau$ are perfect, so $\det b_\tau$
is a unit.  Returning to any other coefficient normalization
or primitive coefficient line changes the rank-six determinant
by a sixth power of an element fixed by bar, and therefore
changes its valuation by an even integer.
Thus the intrinsic $\tau$ determinant norm class has even
valuation at $p$.

HD-PD-01 compares this class with $\delta_B$.
The transfer factor $|N|$ and all coefficient or projector
scales contribute only sixth powers.  This remains true when
$p$ divides $|N|$ or $|G|$.
Thus (PD.9) follows.  Since the norm of an element of the
unramified quadratic extension $K_p/\mathbf Q_p$ has even
valuation, the class $p$ cannot be the signed determinant
class of $B$.

\paragraph{Target existence and arbitrary polarizations.}
The $K$-Hermitian matrix
\[
                    \operatorname{diag}(1,1,1,-1,-1,-p)
 \tag{PD.18}
\]
has signature $(3,3)$ and signed determinant $p$.
A compatible complex structure and a suitable positive
integral multiple of its alternating trace form on an
$\mathcal O_K$-lattice give a polarized abelian variety.
The associated unitary period domain has dimension nine.
The previously checked full-domain centralizer argument
gives $\operatorname{End}^0=K$ for its very general member.

More generally suppose a target with endomorphism algebra
$K$ is $K$-isogenous to a natural source $B$.
Then $\operatorname{End}^0(B)=K$.
Relative to the natural polarization, the pullback of the
target polarization is a positive Rosati-symmetric
endomorphism, hence a positive element of $K^+=\mathbf Q$.
Equivalently, compare the two alternating Hodge pairings:
their ratio is a Hodge endomorphism in $K$, and alternation
makes it fixed by conjugation.  Its rank-six determinant
scale is a norm.  Thus the prescribed target polarization
has the same signed determinant class as the natural source.
This proves the stated isogeny exclusion.  No assertion is
made that a special target with larger endomorphism algebra
has only one polarization determinant type.

\paragraph{Infinitude and an effective construction.}
Dirichlet's theorem gives infinitely many primes
$p\equiv-1\pmod M$
\cite[Section 1, p.~63]{PD:Harrison2009}.
They are inert in $K$ and their norm classes are pairwise
distinct, as seen from the odd valuation at their respective
inert primes.  Compute $M$, enumerate $p=kM-1$ above any
specified bound, and test primality until one is found.
Each primality test is finite, and Dirichlet proves termination.
No GRH, numerical prime estimate, branch-tuple enumeration,
Hodge-locus decision or bound on a minimal original index
enters this algorithm.
\end{proof}

\paragraph{The resulting necessary gate.}
For every $Q$ there are infinitely many determinant types
that have no source with a faithful inducing image of order
at most $Q$.  Therefore any exact-$K$ finite-group primitive-factor
mechanism supplying every determinant type would require
unbounded minimal inducing-image orders across those types.
If a separately proved structural class has a bound $Q$,
HD-PD-02 excludes these types from that class immediately.
The current application does not assume that the
factorwise tensor-decomposable interface in BI19 has been
resolved, or that the bounded structural class is exhaustive.

\paragraph{An explicit warning about proper-field determinant fibers.}
The restriction map on determinant norm classes need not be
injective for an even relative field degree.  For
$K=\mathbf Q(i)$, $E=\mathbf Q(i,\sqrt2)$ and
$F=\mathbf Q(\sqrt2)$,
\[
                 3=\operatorname{Nm}_{E/F}(1+i\sqrt2),
 \tag{PD.19}
\]
whereas $3$ is not a $K$-norm because it has odd valuation at
the inert prime $3$.  The selected primes in HD-PD-02 avoid
this issue: at those primes $E_w=K_p$, so the local obstruction
survives every relevant coefficient-field extension.
